% =====================================================================
%  frontmatter.tex —— 标题页、摘要、译者说明
% =====================================================================
\title{\bfseries 四维规范场论中的纽结不变量\\[6pt]
       \large Knot Invariants from Four-Dimensional Gauge Theory}
\author{Davide Gaiotto$^{1}$ \quad Edward Witten$^{1,2}$\\[4pt]
        \small 中译：学术翻译工作稿}
\date{arXiv:1106.4789v1 [hep-th]，2011 年 6 月 23 日}
\maketitle

\begin{center}
\small
$^1$School of Natural Sciences, Institute for Advanced Study,\\
1 Einstein Drive, Princeton, NJ 08540 USA\\
$^2$Department of Physics, Stanford University, Palo Alto, CA 94305
\end{center}

\vspace{1em}

\begin{abstract}
\noindent
已有论证表明，基于电磁对偶及其他要素，三维中纽结的 Jones 多项式可以通过计数四维中某些规范场方程的解来计算。本文试图不借助任何量子对偶性，直接分析这些方程并计数其解，以验证这一结果。在对这些方程作适当扰动、使其行为更为一般之后，我们能够相当清楚地理解 Jones 多项式是如何浮现的。论证的主要要素，是所讨论的四维规范场方程与二维 Virasoro 代数退化表示的共形块之间的联系。在此过程中，我们对该主题与数学物理中众多新旧课题的关联获得了更好的理解——从 Gaudin 自旋链的 Bethe 拟设，到 BPS 单极子的 $M$ 理论描述，再到 Chern-Simons 规范理论与 Virasoro 共形块之间的关系。
\end{abstract}

\newpage

\vspace{0.5em}

\begin{transnote}
\textbf{关于本译稿。}
\begin{itemize}\setlength{\itemsep}{2pt}
\item \textbf{底本}：英文原稿 D. Gaiotto, E. Witten, \emph{Knot Invariants from Four-Dimensional Gauge Theory}, arXiv:1106.4789v1（120 页，含 29 幅图）。译文以原稿的 \LaTeX{} 源码为底本进行语义级翻译，\emph{全部数学公式逐字保留原文}（仅译出 \verb|\text{...}| 中的说明性文字），因此公式与原稿严格一致；若原稿存在笔误或记号歧义，均以 \verb|% [Tutor Note: ...]| 的形式在源码中就地标注。
\item \textbf{术语}：遵循中文理论物理界约定（规范场论、和乐、反常、共形块、算符乘积展开等）；人名、理论名与数学专名（Jones, Khovanov, Witten, Nahm, oper, thimble 等）保留原文。
\item \textbf{图表}：原 PDF 中的 29 幅图已逐幅裁剪为位图，存于 \verb|figs/| 目录，并在正文对应位置以 \verb|figure| 环境插入；图注中标注了原稿页码，便于对照。
\item \textbf{批注}：正文关键推导处嵌有 \verb|% [Tutor Note: ...]| 形式的学术批注（易错点预警、作者省略的证明、约定差异、原文疑似笔误等），供研读时参考；这些批注不影响编译产物。全部批注另汇编为随稿交付的 \verb|审校批注汇编.md|，按章节顺序排列，便于检索。
\item \textbf{编译}：XeLaTeX（\verb|ctex| 宏包）编译；源码缩进与分段已整理，可直接增删修改。
\end{itemize}
\end{transnote}

\vspace{0.5em}
\noindent\textbf{记号约定速查}\quad
$G$：本文重点研究的“磁”描述中的规范群；$G^\vee$：Chern-Simons 描述中的规范群（与原文 \cite{fiveknots} 的记号相反）；
$\t$：BPS 方程中的实参数，取值于 $\mathbb{RP}^1$；
$b^\pm$：二形式 $b$ 的自对偶/反自对偶投影；
$\d_A$：规范协变外微分；$\Tr$：规范群伴随表示上的迹。
上述约定与原稿一致，阅读时请特别注意 $\t$ 的特殊地位——它并非耦合常数，而是方程族的一个形变参数，$\t=1$ 与一般 $\t$ 的物理差异是全文的一条主线。

\newpage
\tableofcontents
\newpage
